\section{How Trades Are Built: Setup Stacking, Context, and Quality}

\subsection{The corpus does not believe in one-factor trading}

The transcript material is explicit that a trade should be built like a constructor set from multiple components. A graphically obvious level alone is weak. A density alone is weak. An iceberg alone may be interesting but still insufficient. Quality comes from stacked evidence.

\paragraph{Practical implication.} The right way to think about a setup in this project is not ``pattern present or absent'' but rather:

\begin{enumerate}
\item what base ingredients are present,
\item how many of them align,
\item whether any strong counter-signals are present,
\item whether current market conditions are suitable for this specific pattern,
\item whether this pattern appears to be working today for this instrument.
\end{enumerate}

\subsection{Minimum bases for a trade}

The transcript corpus suggests that meaningful trades begin when multiple reasons align. For engineering purposes, a useful translation is:

\begin{itemize}
\item \textbf{one factor:} weak observation, usually not tradable,
\item \textbf{two factors:} hypothesis worth watching,
\item \textbf{three or more factors:} candidate trade,
\item \textbf{three plus strong counter-signals:} candidate rejected,
\item \textbf{three plus confirming sequence plus clean stop:} high-quality candidate.
\end{itemize}

\subsection{Core evidence families}

The recurring evidence families are:

\begin{longtable}{p{0.19\textwidth}p{0.31\textwidth}p{0.42\textwidth}}
\toprule
\textbf{Evidence family} & \textbf{Examples} & \textbf{Why it matters} \\
\midrule
Graph context & prior high, prior low, gap edge, range boundary, local retest, multi-test level & tells us where stops, trapped traders, and attention are likely concentrated \\
Visible liquidity & density, ladder of densities, loaded side of book, empty side of book & defines immediate barriers, magnets, and short-term stop placement logic \\
Hidden / replenishing liquidity & iceberg, repeated refill, unexpectedly persistent level & suggests a stronger participant than the displayed size alone implies \\
Aggressive flow & frequent prints, larger prints, one-sided tape, sudden acceleration & indicates whether the level is being attacked or defended in real time \\
Robot behavior & mechanical reposting, one-tick grinding, steady push & reveals execution logic that other traders can lean on temporarily \\
Session context & opening auction aftermath, news release, lunch slowdown, late-session squeeze & changes both pattern quality and expected follow-through \\
Cross-instrument context & leader instrument, index, currency leg, correlated future & filters whether the local move is supported or fighting the broader tape \\
\bottomrule
\end{longtable}

\subsection{Clean versus unclean setups}

The source corpus often speaks in terms that are partly objective and partly intuitive: a setup can look ``clean'', ``dirty'', ``crooked'', ``late'', ``too empty'', ``already moved'', or ``overfilled with passengers''. A research translation of those judgments is useful.

\paragraph{A clean setup usually has:}
\begin{itemize}
\item clear location in chart context,
\item obvious nearby invalidation,
\item visible support from DOM behavior,
\item tape consistent with the intended direction,
\item acceptable spread and depth,
\item not too much prior exhaustion,
\item reasonable remaining room before the next obstacle.
\end{itemize}

\paragraph{An unclean setup usually has:}
\begin{itemize}
\item late entry after too much movement,
\item unclear stop or stop placed through empty book,
\item conflicting tape and DOM evidence,
\item too many nearby opposing densities,
\item no remaining room,
\item obvious dependence on hope rather than sequence confirmation.
\end{itemize}

\subsection{Day-specific calibration}

One of the strongest project-level ideas supported by the corpus is that the first part of the session should be used as a diagnostic period. The point is not merely to ``wait for setup''; the point is to learn:

\begin{itemize}
\item Are visible densities holding today or being sliced?
\item Are breakouts following through today or failing quickly?
\item Are robot-driven pushes persistent or easy to fade?
\item Is the book deep enough for size, or too thin for confidence?
\item Is the market responsive to tape aggression, or are prints being absorbed with little travel?
\end{itemize}

\researchnote{This naturally suggests an intraday adaptation layer: a running dashboard of setup hit-rates and travel statistics conditioned on the current day and instrument.}

\subsection{Setup taxonomy used in this report}

The formal setup taxonomy used below is:

\begin{enumerate}
\item bounce from visible density,
\item breakout through visible density or chart level,
\item false breakout / failed continuation,
\item iceberg-like support or resistance,
\item absorption and hidden liquidity,
\item robot-driven movement,
\item spring / recoil after one-sided impulse,
\item volatility harvesting and spread capture,
\item leader-follower or guide-instrument alignment,
\item news-response trading,
\item limit-up / limit-down style constraints,
\item context-only filters that modify the interpretation of all of the above.
\end{enumerate}

